%%%%%%%%%%%%%%%%%%%%%%%%%%%%% Define Article %%%%%%%%%%%%%%%%%%%%%%%%%%%%%%%%%%
\documentclass{article}
%%%%%%%%%%%%%%%%%%%%%%%%%%%%%%%%%%%%%%%%%%%%%%%%%%%%%%%%%%%%%%%%%%%%%%%%%%%%%%%

%%%%%%%%%%%%%%%%%%%%%%%%%%%%% Using Packages %%%%%%%%%%%%%%%%%%%%%%%%%%%%%%%%%%
\usepackage{geometry}
\usepackage{graphicx}
\usepackage{amssymb}
\usepackage{amsmath}
\usepackage{amsthm}
\usepackage{empheq}
\usepackage{mdframed}
\usepackage{booktabs}
\usepackage{lipsum}
\usepackage{graphicx}
\usepackage{color}
\usepackage{psfrag}
\usepackage{pgfplots}
\usepackage{bm}
%%%%%%%%%%%%%%%%%%%%%%%%%%%%%%%%%%%%%%%%%%%%%%%%%%%%%%%%%%%%%%%%%%%%%%%%%%%%%%%

% Other Settings

%%%%%%%%%%%%%%%%%%%%%%%%%% Page Setting %%%%%%%%%%%%%%%%%%%%%%%%%%%%%%%%%%%%%%%
\geometry{a4paper}

%%%%%%%%%%%%%%%%%%%%%%%%%% Define some useful colors %%%%%%%%%%%%%%%%%%%%%%%%%%
\definecolor{ocre}{RGB}{243,102,25}
\definecolor{mygray}{RGB}{243,243,244}
\definecolor{deepGreen}{RGB}{26,111,0}
\definecolor{shallowGreen}{RGB}{235,255,255}
\definecolor{deepBlue}{RGB}{61,124,222}
\definecolor{shallowBlue}{RGB}{235,249,255}
%%%%%%%%%%%%%%%%%%%%%%%%%%%%%%%%%%%%%%%%%%%%%%%%%%%%%%%%%%%%%%%%%%%%%%%%%%%%%%%

%%%%%%%%%%%%%%%%%%%%%%%%%% Define an orangebox command %%%%%%%%%%%%%%%%%%%%%%%%
\newcommand\orangebox[1]{\fcolorbox{ocre}{mygray}{\hspace{1em}#1\hspace{1em}}}
%%%%%%%%%%%%%%%%%%%%%%%%%%%%%%%%%%%%%%%%%%%%%%%%%%%%%%%%%%%%%%%%%%%%%%%%%%%%%%%

%%%%%%%%%%%%%%%%%%%%%%%%%%%% English Environments %%%%%%%%%%%%%%%%%%%%%%%%%%%%%
\newtheoremstyle{mytheoremstyle}{3pt}{3pt}{\normalfont}{0cm}{\rmfamily\bfseries}{}{1em}{{\color{black}\thmname{#1}~\thmnumber{#2}}\thmnote{\,--\,#3}}
\newtheoremstyle{myproblemstyle}{3pt}{3pt}{\normalfont}{0cm}{\rmfamily\bfseries}{}{1em}{{\color{black}\thmname{#1}~\thmnumber{#2}}\thmnote{\,--\,#3}}
\theoremstyle{mytheoremstyle}
\newmdtheoremenv[linewidth=1pt,backgroundcolor=shallowGreen,linecolor=deepGreen,leftmargin=0pt,innerleftmargin=20pt,innerrightmargin=20pt,]{theorem}{Theorem}[section]
\theoremstyle{mytheoremstyle}
\newmdtheoremenv[linewidth=1pt,backgroundcolor=shallowBlue,linecolor=deepBlue,leftmargin=0pt,innerleftmargin=20pt,innerrightmargin=20pt,]{definition}{Definition}[section]
\theoremstyle{myproblemstyle}
\newmdtheoremenv[linecolor=black,leftmargin=0pt,innerleftmargin=10pt,innerrightmargin=10pt,]{problem}{Problem}[section]
\theoremstyle{myproblemstyle}
\newmdtheoremenv[linecolor=black,leftmargin=0pt,innerleftmargin=10pt,innerrightmargin=10pt,]{example}{Example}[section]
%%%%%%%%%%%%%%%%%%%%%%%%%%%%%%%%%%%%%%%%%%%%%%%%%%%%%%%%%%%%%%%%%%%%%%%%%%%%%%%

%%%%%%%%%%%%%%%%%%%%%%%%%%%%%%% Plotting Settings %%%%%%%%%%%%%%%%%%%%%%%%%%%%%
\usepgfplotslibrary{colorbrewer}
\pgfplotsset{width=8cm,compat=1.9}
%%%%%%%%%%%%%%%%%%%%%%%%%%%%%%%%%%%%%%%%%%%%%%%%%%%%%%%%%%%%%%%%%%%%%%%%%%%%%%%

%%%%%%%%%%%%%%%%%%%%%%%%%%%%%%% Title & Author %%%%%%%%%%%%%%%%%%%%%%%%%%%%%%%%
\title{Solving Equations}
\author{Alston Yam}
\date{}
\parskip=2pt
\parindent=0pt
%%%%%%%%%%%%%%%%%%%%%%%%%%%%%%%%%%%%%%%%%%%%%%%%%%%%%%%%%%%%%%%%%%%%%%%%%%%%%%%

\begin{document}
    \maketitle

    \section*{Learning Goals}
    \begin{itemize}
        \item Understand what a simultaneous equation represents
        \item Solving systems with one variable
        \item Solve systems with substitution and elimination
        \item Check if your solution makes sense    
    \end{itemize}

    \section{Why Equations?}
    You might be thinking about $x$ and $y$. These are simply just a \textit{place holder} for a number you don't know yet. We call it a \textbf{variable}.

    \begin{example}
        I have a number, such that when I multiply it by $4$ and then add $2$, I get $18$. What is my number?
    \end{example}

    It's very easy to see that my number is $4$. But we will put it in equation terms:
    \[4x + 2 = 18\]

    I'm just using $x$ as a place holder for the number that I don't know yet. 

    \begin{example}
        I have two numbers. They add to $10$, and three times the first number take away the second number is $2$. What are the two numbers?
    \end{example}
    
    Now we have two numbers, we can use $x$ and $y$.
    \begin{align*}
        x + y &= 10 \\
        3x - y &= 2
    \end{align*}

    \begin{problem}
        Sarah has some apples and oranges. She has 3 more apples than oranges. If she has 15 pieces of fruit altogether, how many apples and oranges does she have?
    \end{problem}

    \begin{problem}
        A movie ticket costs \$8 and popcorn costs \$5. If Marcus spent \$46 on 7 items total (tickets and popcorns), how many of each did he buy?
    \end{problem}
    

    \section{One Variable — The Basics}

    You need to treat each equation like a balance. Whatever you do to one side, you \textbf{must} do it to the other side too, so the balance don't tip.

    We can look at the previous example:
    \[4x + 2 - 2 = 18 - 2 \implies 4x = 16 \implies x = 4\]

    Kind of similar to working backwards, right?

    \begin{problem}
        Solve for $x$: $2x - 5 = 11$
    \end{problem}

    \begin{problem}
        Solve for $x$: $\frac{x}{3} + 4 = 7$
    \end{problem}

    \begin{problem}
        Solve for $x$: $5(x - 2) = 30$
    \end{problem}

    \begin{problem}
        Solve for $x$: $7x + 3 = 2x + 18$
    \end{problem}


    \section{Two Variables}
    Lets look at our example from before. 
    \begin{align*}
        x + y &= 10 \\
        3x - y &= 2
    \end{align*}

    \subsection{Substitution}
    We are going to \textit{isolate} the variables.

    \[x + y = 10 \implies y = 10 - x\]

    and 
    \[3x - y = 2 \implies 3x = 2 + y \implies 3x - 2 = y \implies y = 3x - 2\]

    But wait, we now have $y = \text{something}$ in two ways. Since both are representing the same thing, lets put them together:

    \[10 - x = 3x - 2\]

    We can solve this now for $x$, and we get $x = 3$. WAIT! We aren't done yet. We only found $x$, but it's very important you also find $y$.

    To do that, we can just change $x$ into $3$ in one of the equations we had.

    \[3 + y = 10 \implies y = 7\]

    \subsection{Elimination}

    Lets take the same equations again:
    \begin{align*}
        x + y &= 10 \\
        3x - y &= 2
    \end{align*}

    In the end, the thing on the left side $x+y$ is just another number, and we know that number is $10$. Same thing with $3x - y$. So actually we can just add both of these equations together.
    \begin{align*}
        x + y + (3x - y) &= 10 + 2 \\
        4x &= 12 \\
        x &= 3
    \end{align*}

    Same thing as before, now we have $x$, we substitute it back and get $y = 7$.

    \subsection{Presenting \& Checking Your Solution}
    After you've solved the equation and you have your $x$ and $y$, write it like this: $(x, y) = (3, 7)$.

    To check your solution, you need to substitute your answers back.
    \begin{align*}
        3 + 7 &= 10 \\
        3 \times 3 - 7 &= 2
    \end{align*}

    Which all works out!

    \subsection{Practice Problems}
    \begin{problem}
        Solve the system using substitution or elimination:
        \begin{align*}
            x + 2y &= 7 \\
            2x - y &= 4
        \end{align*}
    \end{problem}

    \begin{problem}
        Solve the system using substitution or elimination:
        \begin{align*}
            3x + y &= 11 \\
            x - y &= -3
        \end{align*}
    \end{problem}

    \begin{problem}
        Solve the system using substitution or elimination:
        \begin{align*}
            2x + 3y &= 12 \\
            4x - y &= 5
        \end{align*}
    \end{problem}

    \pagebreak

    \section{Word Problems}
    \begin{problem}
        A farmer has a total of 50 animals, consisting of chickens and cows. If there are 20 more chickens than cows, how many of each type of animal does the farmer have?
    \end{problem}

    \begin{problem}
        A bookstore sold a total of 120 books in one day. If the number of fiction books sold was twice the number of non-fiction books sold, how many of each type were sold?
    \end{problem}

    \begin{problem}
        A class has a total of 30 students. If there are 5 more girls than boys in the class, how many boys and girls are there?
    \end{problem}

    \begin{problem}
        A train leaves a station traveling at a speed of 60 km/h. Another train leaves the same station 30 minutes later, traveling at a speed of 90 km/h. How far from the station will the two trains meet?
    \end{problem}

    \begin{problem}
        A rectangle has a length that is twice its width. If the perimeter of the rectangle is 48 cm, what are the dimensions of the rectangle?
    \end{problem}

    \begin{problem}
        The sum of three consecutive integers is 72. What are the integers?
    \end{problem}

    \begin{problem}
        A car rental company charges a flat fee of \$50 plus \$0.20 per mile driven. If a customer pays \$70, how many miles did they drive?
    \end{problem}

    \begin{problem}
        A farmer has a total of 35 heads of chickens and rabbits. If the total number of feet is 94, how many chickens and how many rabbits does the farmer have?
    \end{problem}
    
    
\end{document}